% ============================================================================
\clearpage
\section*{v4.9.0-r78 Source-Only: Covariance of the OP5 Comparison System at Value Level and the Stabilizer of the Gauge}
\addcontentsline{toc}{section}{v4.9.0-r78 Source-Only: Covariance of the OP5 Comparison System at Value Level and the Stabilizer of the Gauge}
\label{sec:v490-r78-value-covariance}
% ============================================================================

\paragraph{Purpose.}
Layer v4.9.0-r77 showed that the construction \(R\) of the OP5 comparison system is covariant at the level of its first-order
symbol (H64, \S\ref{sec:v490-r77-high-precision-covariance}).  On 4 November 2026 the author asked to check the remaining
building blocks: \(\Omega\) and \(\theta\), the four hidden conditions and the gauge \(\Psi^\dagger K_2\Psi=S\).  \(\Omega\) is
derived from \(R\) alone, the four conditions from \(R\) together with the gauge.  Hence it suffices to establish two things:
\begin{itemize}[leftmargin=2em]
\item the covariance of \(R\) at value level, with all lower-order terms;
\item the stabilizer of the gauge.
\end{itemize}
This layer records one pre-registered test that is invalid by its criterion (H65) and one confirmed test (H66).  Everything
concerns the comparison model of \S\ref{sec:v410-state-mediator} only.  Earlier layers are not rewritten.

\noindent\textbf{Source-only status.}
This v4.9.0-r78 layer is source-only.  It does not build a PDF, does not assign a new DOI,
and does not authorize a public release.  The v4.9.0 Version DOI
\href{https://doi.org/10.5281/zenodo.23534954}{10.5281/zenodo.23534954} and the Concept DOI
\href{https://doi.org/10.5281/zenodo.19210728}{10.5281/zenodo.19210728} are unchanged by it.

\subsection*{v4.9.0-r78 results}
\label{subsec:v490-r78-results}

\paragraph{Record v4.9.0-r78.1 (H65 invalid: stabilizer of the gauge as predicted, value level below the control bound).}
Pre-registered.
\begin{itemize}[leftmargin=2em]
\item \emph{Part (b), the stabilizer of the gauge.}  The generators \(X\in\mathfrak{so}(1,3)\), in the spin representation, that
preserve \(\Psi^\dagger K_2\Psi\) are those with \(K_2X+X^\dagger K_2=0\).  They span a three-dimensional subalgebra: the
rotation about \(y\) and the boosts along \(x\) and \(z\).  This is \(\mathfrak{so}(1,2)\), the stabilizer of \(e_y\), as
predicted.  The gauge is a non-covariant gauge fixing that singles out one spatial direction.
\item \emph{Part (a), value level.}  Run 1 was invalid because the row transformation was taken from the symbol, which does not
determine it off the solution set (positive control \(0.24\)).  A dated addendum preceded run 2, which is not blind and fits
the row transformation from values.  Run 2 gave residuals \(\le1.1\times10^{-10}\), but its positive control
(\(5.6\times10^{-11}\)) missed the pre-registered bound \(10^{-12}\), which lay below the numerical floor.  Part (a) is
therefore invalid in both runs, and so is H65 as a whole.  The bound was not relaxed after the fact.
\end{itemize}

\paragraph{Theoremlet v4.9.0-r78.2 (H66 confirmed: \(R\) is covariant at value level).}
Pre-registered, not blind (the values of H65 were known); the control bound is \(10^{-9}\).
\begin{itemize}[leftmargin=2em]
\item \emph{Setup.}  For each fixed pair \((M,S)\), with \(M\in GL(4)\) and \(S\in\mathrm{Spin}(1,3)\), the row
transformation \(C\) is fitted from 80 random data sets off the solution set: \(\Psi\), its jets, and the background
\(A,\partial A,h,\partial h\).  It is tested on 20 further data sets.  Five pairs were used: a pure rotation, a pure boost, a
pure change of coordinates, and two mixed pairs.
\item \emph{Result.}  All relative residuals are \(\le2.2\times10^{-11}\), and \(C\) is block diagonal on \(F,G,dg\) to
\(5\times10^{-11}\).  The positive control gives \(\le3.2\times10^{-13}\).  The negative control, with the background held
fixed, gives \(57\) for \(F\) and \(1.2\times10^3\) for \(G\).
\item \emph{Reading.}  \(R\) is covariant with all lower-order terms.  Hence \(\Omega\), which is derived from \(R\) alone, is
covariant.  The four hidden conditions are derived from \(R\) together with the gauge, so they are covariant at most under
\(\mathfrak{so}(1,2)\).
\end{itemize}

\paragraph{Proof and measurement sources.}
r78.1 by \texttt{k165\_kovarianz\_rest.py} (run 1) and \texttt{k165b\_werte.py} (run 2); r78.2 by
\texttt{k166\_kovarianz\_werte.py}.  The notes \texttt{Vorab\_\allowbreak Notiz\_\allowbreak H65}, \texttt{H66} and the
result notes lie in \texttt{notes/\allowbreak research/\allowbreak predictions/}.  \(\square\)

\paragraph{Consequence.}
The comparison system \(R\) and its integrability condition \(\Omega\) are covariant.  Covariance is broken in two places:
\begin{itemize}[leftmargin=2em]
\item by the prescribed background \((\Psi\text{ at the germ},A,h)\);
\item by the gauge, which fixes the \(y\)-component of a current and leaves \(\mathfrak{so}(1,2)\).
\end{itemize}
Whether the number of the four hidden conditions and their support on the cone survive a change of gauge is open.  So is a
native dynamics for the background.  OP5 is not closed.

\subsection*{v4.9.0-r78 status in the three frames}
\label{subsec:v490-r78-frames}

\begin{itemize}[leftmargin=2em]
\item All statements concern the fixed OP5 comparison model in one chart, with one coframe reconstructed from \(\Psi\).  No
frame map M, T or P acts on them.  The Lorentz transformations here act inside the model and are not the photon frame P.
\end{itemize}
No statement of this layer concerns colour or charge.

\subsection*{v4.9.0-r78 register entries}
\label{subsec:v490-r78-register}

\begin{center}
\small
\begin{tabular}{@{}>{\raggedright\arraybackslash}p{0.34\textwidth}>{\raggedright\arraybackslash}p{0.58\textwidth}@{}}
\toprule
statement & status after v4.9.0-r78\\
\midrule
covariance of \(R\) at value level & confirmed (H66, not blind); H65(a) invalid by its control bound\\
covariance of \(\Omega\) & follows, since \(\Omega\) is derived from \(R\) alone\\
gauge \(\Psi^\dagger K_2\Psi=S\) & preserved exactly by \(\mathfrak{so}(1,2)\), the stabilizer of \(e_y\) (H65(b))\\
four hidden conditions & covariant at most under \(\mathfrak{so}(1,2)\); gauge independence open\\
next step & gauge independence of the four conditions; native dynamics of the background\\
\bottomrule
\end{tabular}
\end{center}

\subsection*{v4.9.0-r78 ledger and guard update}
\label{subsec:v490-r78-ledger-guard}

\[
R_{\rm chk}^{v490\text{-}r78}=(1,1,0)
\]
for (covariance of \(R\) and \(\Omega\) established; H65 recorded as invalid; OP5 not closed).

The layer spends no guard:
\[
R_{\rm glob}^{(v490\text{-}r78,+)}=(0,1,1,1,1),
\qquad
\Gamma_{\rm guard}^{v490\text{-}r78}=(1,1,1,1,1),
\qquad
\beta_4=1/256\ \text{unspent}.
\]
It does not imply exact finite \(C5\) zero closure, exact all-readout terminality,
\(C6/T2\), a global photon-clock/RH theorem, or global all-readout closure.
